% Auto-generated by extract-marking.py -- do not edit by hand unless you know what to do.
% Works for BOTH \documentclass[...,type=solutions,...] and
% [...,type=marking,...] compiles, chosen at compile time via
% \OlympJuryPick / \OlympJuryBeginTable (see olympiad-marking.sty).
% \eqref{n} falls back to a plain '(n)' whenever the referenced
% equation isn't actually typeset in the current document.

\OlympJuryBeginTable{|p{0.862\linewidth}|>{\centering\arraybackslash}p{0.098\linewidth}|}{|p{0.655\linewidth}|>{\centering\arraybackslash}p{0.085\linewidth}|>{\centering\arraybackslash}p{0.06\linewidth}|>{\centering\arraybackslash}p{0.06\linewidth}|}
\hline
\OlympJuryPick{\multicolumn{1}{|c|}{\textbf{Содержание}} & \textbf{Баллы} \\ \hline}{\multicolumn{1}{|c|}{\textbf{Содержание}} & \textbf{Баллы} & \textbf{MP} & \textbf{A} \\ \hline}
\OlympJuryPick{\quad\textit{Критерии за альтернативный метод эквивалентны критериям за описание фазовых диаграмм \textbf{(ФД)}. Общий балл ставится по наибольшему между двумя методами внутри каждого из двух частей задачи. Переноса ошибки на численные расчёты нет.} &  \\ \hline}{\quad\textit{Критерии за альтернативный метод эквивалентны критериям за описание фазовых диаграмм \textbf{(ФД)}. Общий балл ставится по наибольшему между двумя методами внутри каждого из двух частей задачи. Переноса ошибки на численные расчёты нет.} &  & &  \\ \hline}
\OlympJuryPick{Формула \eqref{1}: $T_0 = 2\pi\sqrt{LC}.$ ; принимается $\omega=\sqrt{LC}$ & 0.2 \\ \hline}{Формула \eqref{1}: $T_0 = 2\pi\sqrt{LC}.$ ; принимается $\omega=\sqrt{LC}$ & 0.2 & &  \\ \hline}
\OlympJuryPick{Формула \eqref{2}: $\frac T{T_0} = \frac14$ & 0.1 \\ \hline}{Формула \eqref{2}: $\frac T{T_0} = \frac14$ & 0.1 & &  \\ \hline}
\OlympJuryPick{Формула \eqref{3}: $L\cdot\frac{dI}{dt} + \frac qC = U(t),$ ; принимается $U(t)\Leftrightarrow U_0\Leftrightarrow0$ в правой части & 0.1 \\ \hline}{Формула \eqref{3}: $L\cdot\frac{dI}{dt} + \frac qC = U(t),$ ; принимается $U(t)\Leftrightarrow U_0\Leftrightarrow0$ в правой части & 0.1 & &  \\ \hline}
\OlympJuryPick{Формула \eqref{4}: $\Delta I = \frac{U_0\tau}L.$ & 0.2 \\ \hline}{Формула \eqref{4}: $\Delta I = \frac{U_0\tau}L.$ & 0.2 & &  \\ \hline}
\OlympJuryPick{\textbf{Метод фазовых диаграмм:} (следующие 5 пунктов эквивалентны \eqref{6}--\eqref{7}) & \textbf{?} \\ \hline}{\textbf{Метод фазовых диаграмм:} (следующие 5 пунктов эквивалентны \eqref{6}--\eqref{7}) & \textbf{?} & &  \\ \hline}
\OlympJuryPick{\textbf{(ФД)} Движение фазового состояния по окружности \textbf{и} по часовой стрелке & 0.1 \\ \hline}{\textbf{(ФД)} Движение фазового состояния по окружности \textbf{и} по часовой стрелке & 0.1 & &  \\ \hline}
\OlympJuryPick{\textbf{(ФД)} Нормализация $I/\sqrt{LC}$ или эквивалент & 0.1 \\ \hline}{\textbf{(ФД)} Нормализация $I/\sqrt{LC}$ или эквивалент & 0.1 & &  \\ \hline}
\OlympJuryPick{\textbf{(ФД)} Каждый импульс --- это вертикальный вектор перемещения точки & 0.1 \\ \hline}{\textbf{(ФД)} Каждый импульс --- это вертикальный вектор перемещения точки & 0.1 & &  \\ \hline}
\OlympJuryPick{\textbf{(ФД)} Ввиду \eqref{2} происходит поворот на $\pi/2$ между каждым импульсом & 0.1 \\ \hline}{\textbf{(ФД)} Ввиду \eqref{2} происходит поворот на $\pi/2$ между каждым импульсом & 0.1 & &  \\ \hline}
\OlympJuryPick{\textbf{(ФД)} Верно нарисовано движение точки (как на Рис. \figref{2}) & 0.3 \\ \hline}{\textbf{(ФД)} Верно нарисовано движение точки (как на Рис. \figref{2}) & 0.3 & &  \\ \hline}
\OlympJuryPick{Ответ к \textbf{(a)} \eqref{5}: $E=0.$ ; без обоснования не принимается & 0.4 \\ \hline}{Ответ к \textbf{(a)} \eqref{5}: $E=0.$ ; без обоснования не принимается & 0.4 & &  \\ \hline}
\OlympJuryPick{\quad\textit{\textbf{Альтернативный метод:} (вместо \textbf{ФД})} &  \\ \hline}{\quad\textit{\textbf{Альтернативный метод:} (вместо \textbf{ФД})} &  & &  \\ \hline}
\OlympJuryPick{\quad\textit{Формула} \eqref{6}: $q(t)=\frac{\Delta I}{\omega}\sin(\omega t).$ & \textit{0.2} \\ \hline}{\quad\textit{Формула} \eqref{6}: $q(t)=\frac{\Delta I}{\omega}\sin(\omega t).$ & \textit{0.2} & &  \\ \hline}
\OlympJuryPick{}{\pagebreak}
\OlympJuryPick{\quad\textit{Формула} \eqref{7}: $q(t)=\frac{\Delta I}\omega\ab(\sin\omega(t-T)+\cos\omega(t-T))=\frac{\sqrt2\Delta I}\omega\sin\ab(\omega (t-T)+\frac\pi4).$ & \textit{0.5} \\ \hline}{\quad\textit{Формула} \eqref{7}: $q(t)=\frac{\Delta I}\omega\ab(\sin\omega(t-T)+\cos\omega(t-T))=\frac{\sqrt2\Delta I}\omega\sin\ab(\omega (t-T)+\frac\pi4).$ & \textit{0.5} & &  \\ \hline}
\OlympJuryPick{Формула \eqref{8}: $q = \sqrt2\cdot\Delta I\sqrt{LC}.$ ; принимается факт, что внешняя окружность в $\sqrt2$ раза больше внутренней, либо факт, что энергия после второго импульса в $2^{1/4}$ раз больше & 0.3 \\ \hline}{Формула \eqref{8}: $q = \sqrt2\cdot\Delta I\sqrt{LC}.$ ; принимается факт, что внешняя окружность в $\sqrt2$ раза больше внутренней, либо факт, что энергия после второго импульса в $2^{1/4}$ раз больше & 0.3 & &  \\ \hline}
\OlympJuryPick{\quad\textit{Забыт множитель $\sqrt{LC}$} & \textit{-0.2} \\ \hline}{\quad\textit{Забыт множитель $\sqrt{LC}$} & \textit{-0.2} & &  \\ \hline}
\OlympJuryPick{Ответ к \textbf{(a)} \eqref{9}: $q = U_0\tau\sqrt\frac{2C}{L}=2.09\pu{мКл}.$ & 0.3 \\ \hline}{Ответ к \textbf{(a)} \eqref{9}: $q = U_0\tau\sqrt\frac{2C}{L}=2.09\pu{мКл}.$ & 0.3 & &  \\ \hline}
\OlympJuryPick{Численный ответ в \eqref{9} & 0.1 \\ \hline}{Численный ответ в \eqref{9} & 0.1 & &  \\ \hline}
\OlympJuryPick{\textbf{Метод фазовых диаграмм:} (следующий 1 пункт эквивалентен \eqref{11}) & \textbf{?} \\ \hline}{\textbf{Метод фазовых диаграмм:} (следующий 1 пункт эквивалентен \eqref{11}) & \textbf{?} & &  \\ \hline}
\OlympJuryPick{\textbf{(ФД)} Указано, что $I/\sqrt{LC}$ должен опираться хордой под углом 90 градусов на окружность $q_*$ & 0.3 \\ \hline}{\textbf{(ФД)} Указано, что $I/\sqrt{LC}$ должен опираться хордой под углом 90 градусов на окружность $q_*$ & 0.3 & &  \\ \hline}
\OlympJuryPick{Ответ к \textbf{(b)} \eqref{10}: $t_* = \frac T2=125\pu{мс}.$ ; за отсутствие численного значения штрафа нет & 0.3 \\ \hline}{Ответ к \textbf{(b)} \eqref{10}: $t_* = \frac T2=125\pu{мс}.$ ; за отсутствие численного значения штрафа нет & 0.3 & &  \\ \hline}
\OlympJuryPick{\quad\textit{\textbf{Альтернативный метод:} вместо \textbf{ФД}} &  \\ \hline}{\quad\textit{\textbf{Альтернативный метод:} вместо \textbf{ФД}} &  & &  \\ \hline}
\OlympJuryPick{\quad\textit{Формула} \eqref{11}: $-\omega q_*\sin\omega t_* = -\frac{\Delta I}{2},$ & \textit{0.3} \\ \hline}{\quad\textit{Формула} \eqref{11}: $-\omega q_*\sin\omega t_* = -\frac{\Delta I}{2},$ & \textit{0.3} & &  \\ \hline}
\OlympJuryPick{Ответ к \textbf{(b)} \eqref{12}: $q_* = \frac1{\sqrt2}\cdot\Delta I\sqrt{LC}=U_0\tau\sqrt\frac C{2L}=1.04\pu{мКл}.$ ; балл не ставится если забыт множитель $\sqrt{LC}$ & 0.4 \\ \hline}{Ответ к \textbf{(b)} \eqref{12}: $q_* = \frac1{\sqrt2}\cdot\Delta I\sqrt{LC}=U_0\tau\sqrt\frac C{2L}=1.04\pu{мКл}.$ ; балл не ставится если забыт множитель $\sqrt{LC}$ & 0.4 & &  \\ \hline}
\OlympJuryPick{Численный ответ в \eqref{12} & 0.1 \\ \hline}{Численный ответ в \eqref{12} & 0.1 & &  \\ \hline}
\OlympJuryPick{\textbf{Итого} & \textbf{3.5} \\ \hline}{\textbf{Итого} & \textbf{3.5} &  &  \\ \hline}
\end{tabularx}
